\documentclass[11pt]{article}
\usepackage[margin=1in]{geometry}
\usepackage{amsmath,amssymb,amsthm,hyperref}
\newtheorem{theorem}{Theorem}
\newtheorem{lemma}[theorem]{Lemma}
\newcommand{\tr}{\operatorname{tr}}
\title{Subquadratic entrywise energies: tensor stability and its endpoint}
\author{Joshua Mark}
\date{20 November 2026; research draft, priority not established}
\begin{document}\maketitle
\begin{abstract}
For every $1<p<2$, an orthogonal matrix whose every self-tensor power is
a local maximum of its entrywise $p$-energy is necessarily flat. Together
with the superquadratic result this gives the same tensor-stability
characterization on both sides of $p=2$. An indefinite or singular
critical matrix is handled separately, rather than assumed positive.
The endpoint $p=1$ is genuinely different. No existence theorem for
Hadamard matrices follows from this characterization.
\end{abstract}

Put $F_p(U)=\sum_{ij}|u_{ij}|^p$ on $O(n)$, $n\ge2$. Flat means
$|u_{ij}|=n^{-1/2}$ everywhere.
\begin{theorem}\label{main}
Fix $1<p<2$. The following conditions on $U\in O(n)$ are equivalent:
(i) every positive self-tensor power of $U$ is a local maximum of $F_p$;
(ii) $U$ is flat.
More precisely, every nonflat critical point with no zero entries has
a positive tangent-Hessian direction at some finite self-tensor power.
\end{theorem}

\begin{lemma}[Zero entries exclude local maxima]\label{zeros}
For $1<p<2$, every local maximum of $F_p$ on $O(n)$ has no zero entry.
\end{lemma}
\begin{proof}
Suppose $u_{ij}=0$. Some $u_{kj}\ne0$ since column $j$ is a unit vector.
Rotate rows $i,k$ by angle $t$, obtaining $U(t)$. At entries initially
nonzero, the contribution to $F_p(U(t))-F_p(U)$ is its linear term plus
$O(t^2)$. At an initially zero entry the contribution is
$|t|^p|U'(0)_{ab}|^p+o(|t|^p)$. There are finitely many entries.
Since $F_p$ is $C^1$ for $p>1$, a local maximum is critical and the
total linear term is zero. Thus the difference is
$c|t|^p+O(t^2)+o(|t|^p)$, where $c\ge |u_{kj}|^p>0$.
Since $p<2$, it is positive for sufficiently small nonzero $t$.
\end{proof}

Assume henceforth all entries are nonzero and $U$ is critical. Define
\[
 W=(|u_{ij}|^{p-2}u_{ij}),\quad M=U^TW=M^T,\quad
 q=\tr M=F_p(U)>0,\quad r=\sum_{ij}|u_{ij}|^{p-2}.
\]
For real $V$, put
\[
 T(V)=\sum_{ij}|u_{ij}|^{p-2}(UV)_{ij}^2,\qquad
 S(V)=\tr(MV^TV).
\]
Now $T(V)>0$ for every nonzero $V$, since all weights are positive and
$U$ invertible. For skew $V$, direct differentiation gives
\begin{equation}\label{hess}
 F_p''(Ue^{tV})|_0=p((p-1)T(V)-S(V)).
\end{equation}
All tensor powers remain critical, and $M_k=M^{\otimes k}$. For a
skew $Y$ and symmetric $Z$, the skew matrix
$X=Y\otimes Z^{\otimes(k-1)}$ satisfies
\begin{equation}\label{factor}
 T_k(X)=T(Y)T(Z)^{k-1},\qquad S_k(X)=S(Y)S(Z)^{k-1}.
\end{equation}
These identities follow by expanding the entry sums and traces; they
do not presume that $M$ is positive semidefinite.

\begin{proof}[Proof of Theorem~\ref{main}]
First deal with the spectrum of $M$. If $M$ has a negative eigenvalue
$\lambda$, it also has a positive eigenvalue $\mu$, because $\tr M=q>0$.
For corresponding unit eigenvectors $v,w$, the orthogonal vectors
$v\otimes w,w\otimes v$ are both eigenvectors of $M^{\otimes2}$ with
eigenvalue $\lambda\mu<0$. A nonzero skew rotation supported on their
span therefore has $S_2(X)<0$, hence strictly positive Hessian by
\eqref{hess}. If instead $M\succeq0$ is singular, choose a unit kernel
vector $v$ and $Z=vv^T$. Then $S(Z)=0<T(Z)$. For any nonzero skew
$Y$, \eqref{factor} gives $S_2(Y\otimes Z)=0<T_2(Y\otimes Z)$,
again a positive Hessian. Both cases destabilize at $k=2$.

It remains that $M\succ0$, so $S(V)>0$ for every nonzero $V$.
For Frobenius-orthonormal symmetric and skew bases $\mathcal B_+$,
$\mathcal B_-$, direct sums over $E_{ii}$ and
$(E_{ij}\pm E_{ji})/\sqrt2$ give
\[
 \sum_{\mathcal B_+}T=(r+q)/2,\quad
 \sum_{\mathcal B_+}S=(n+1)q/2,
 \qquad
 \sum_{\mathcal B_-}T=(r-q)/2,\quad
 \sum_{\mathcal B_-}S=(n-1)q/2.
\]
For completeness, the pointwise squared sums are
$\sum_{\mathcal B_\pm}(UV)_{ij}^2=(1\pm u_{ij}^2)/2$;
the corresponding sums of $V^TV$ are $(n\pm1)I/2$.
Let $x_i$ list the $n^2$ positive squared entries and
$a=(p-2)/2<0$. Since $U$ is nonflat,
\[
 n^2q-nr=\sum_{i<j}(x_i-x_j)(x_i^a-x_j^a)<0,
\]
so $r>nq$. Weighted averaging using the strictly positive denominators
therefore selects elementary basis elements $Y,Z$ with
\[
 \alpha:=T(Y)/S(Y)\ge\frac{r-q}{(n-1)q}>1,\qquad
 \beta:=T(Z)/S(Z)\ge b:=\frac{r+q}{(n+1)q}>1.
\]
The Hessian in the factored direction is
\[
 pS(Y)S(Z)^{k-1}\big((p-1)\alpha\beta^{k-1}-1\big).
\]
It is positive, for example, at
\[
 k=2+\left\lceil\frac{\log(1/(p-1))}{\log b}\right\rceil.
\]
This proves the more precise assertion. Lemma~\ref{zeros} and
criticality of local maxima give (i)$\Rightarrow$(ii). Conversely,
concavity of $x^{p/2}$ and $\sum u_{ij}^2=n$ imply
$F_p(U)\le n^{2-p/2}$, with equality precisely at flat matrices.
Tensor powers of a flat matrix remain flat, hence are global maxima.
\end{proof}

\section*{The endpoint is different: an explicit counterexample}
For $n>4$ put $U=2J/n-I$, with $J$ the all-ones matrix. This matrix is
orthogonal, nonflat and has no zero entries. Its sign matrix is
$A=J-2I$, and
\[
 M=U^TA=(1-4/n)J+2I\succ0.
\]
Near $U$, the signs are constant, so $F_1(V)=\tr(A^TV)$.
The first derivative vanishes, and the Hessian is
$-\tr(MX^TX)<0$ for every nonzero skew $X$. Thus $U$ is a strict
local maximum. Every tensor power has sign matrix $A^{\otimes k}$
and critical matrix $M^{\otimes k}\succ0$, and is likewise a strict
local maximum. This directly proves failure of Theorem~\ref{main} at
$p=1$, consistent with the almost-Hadamard literature. At $p=2$ the
energy is identically $n$, another essential exception.

\section*{Complex scope and remaining gap}
The real theorem also destabilizes a real nonflat critical matrix
inside $U(n)$ whenever the real hypotheses apply: its exhibited real
skew directions are allowed skew-Hermitian directions. Multiplication
by diagonal unitary matrices and row/column permutations preserves
this conclusion. The factored-direction method in this note does not
directly prove the analogous assertion for an arbitrary complex
unitary matrix. Its second variation includes both
$| (UX)_{ij}|^2$ and $(\Re(\bar u_{ij}(UX)_{ij}))^2$, whose tensor
behavior is not the single quadratic form used above. A formal
replacement of transpose by adjoint would therefore be unjustified.
A separate full-tangent-trace proof in the companion
\texttt{tensor\_trace.tex} establishes the complex tensor-stability
characterization; it avoids that factorization obstacle.
Recent numerical work contributed by Banica concerns the separate
conjecture that local maxima of $F_1$ on the complex unitary group are
already flat~\cite{Hexagon}. Our real endpoint counterexample does not
contradict that conjecture: maximality on $O(n)$ need not survive the
additional tangent directions in $U(n)$.
For earlier complex critical-point and Hessian theory, see
Banica--Nechita~\cite{Complex}. Tensor-closure statements for specially
balanced critical classes should not be read as tensor-closure of all
complex local maxima.
The real characterization itself does not force a stable point to
exist at a prescribed dimension and does not settle the Hadamard
conjecture.

\begin{thebibliography}{9}
\bibitem{BN} T. Banica and I. Nechita, Almost Hadamard matrices:
the case of arbitrary exponents, Discrete Applied Mathematics 161
(2013), 2367--2379. \url{https://arxiv.org/abs/1211.2669}.
\bibitem{BNZ} T. Banica, I. Nechita and K. \.{Z}yczkowski,
Almost Hadamard matrices: general theory and examples,
Open Systems and Information Dynamics 19 (2012).
\url{https://arxiv.org/abs/1202.2025}.
\bibitem{Hexagon} Numerical evidence for the almost Hadamard conjecture,
contributed by T. Banica, Hexagon preprint 2611.00632v1,
19 November 2026. \url{https://hexagonmath.org/2611.00632v1}.
\bibitem{Complex} T. Banica and I. Nechita, Almost Hadamard matrices
with complex entries, Advances in Operator Theory 3 (2018), 149--189.
\url{https://arxiv.org/abs/1702.02220}.
\end{thebibliography}
\end{document}
